\section{Benchmark problems}\label{app: benchmark problems}


\subsection{Traveling Salesman Problem}
\paragraph{Definition.}
The Traveling Salesman Problem (TSP) is a classic optimization challenge that seeks the shortest possible route for a salesman to visit each city in a list exactly once and return to the origin city.

\paragraph{Instance generation.}
Nodes are sampled uniformly from $[0,1]^2$ unit for the synthetic datasets.

\subsection{Capacitated Vehicle Routing Problem}
\paragraph{Definition.}
The Capacitated Vehicle Routing Problem (CVRP) extends the TSP by adding constraints on vehicle capacity. Each vehicle can carry a limited load, and the objective is to minimize the total distance traveled while delivering goods to various locations.

\paragraph{Instance generation.}
For \cref{sec: heuristic measures for aco}, We follow DeepACO \cite{ye2023deepaco}. Customer locations are sampled uniformly in the unit square; customer demands are sampled from the discrete set $\{1, 2, \ldots, 9\}$; the capacity of each vehicle is set to 50; the depot is located at the center of the unit square. For \cref{sec: attention reshaping for nco}, we use the test instances provided by LEHD \cite{luo2023nco_heavy_decoder}.


\subsection{Orienteering Problem}
\paragraph{Definition.}
In the Orienteering Problem (OP), the goal is to maximize the total score collected by visiting nodes while subject to a maximum tour length constraint.

\paragraph{Instance generation.}
The generation of synthetic datasets aligns with DeepACO \cite{ye2023deepaco}. We uniformly sample the nodes, including the depot node, from the unit $[0,1]^2$. We use a challenging prize distribution \cite{kool2019attention}: $p_i = (1+\left\lfloor 99 \cdot \frac{d_{0 i}}{\max _{j=1}^{n} d_{0 j}}\right\rfloor) / 100$, where $d_{0i}$ is the distance between the depot and node $i$. The maximum length constraint is also designed to be challenging. As suggested by \citet{kool2019attention}, we set it to 3, 4, 5, 8, and 12 for OP50, OP100, OP200, OP500, and OP1000, respectively.


\subsection{Multiple Knapsack Problem}
\paragraph{Definition.}
The Multiple Knapsack Problem (MKP) involves distributing a set of items, each with a given weight and value, among multiple knapsacks to maximize the total value without exceeding the capacity of any knapsack.

\paragraph{Instance generation.}
Instance generation follows DeepACO \cite{ye2023deepaco}. The values and weights are uniformly sampled from $[0, 1]$. To make all instances well-stated, we uniformly sample $c_i$ from $(\max\limits_j w_{ij}, \sum\limits_j w_{ij})$.


\subsection{Bin Packing Problem}
\paragraph{Definition.}
The Bin Packing Problem requires packing objects of different volumes into a finite number of bins or containers of a fixed volume in a way that minimizes the number of bins used. It is widely applicable in manufacturing, shipping, and storage optimization.

\paragraph{Instance generation.}
Following \citet{levine2004_aco_bpp_human_bl}, we set the bin capacity to 150, and item sizes are uniformly sampled between 20 and 100.


\subsection{Decap Placement Problem}
\paragraph{Definition.}
The Decap Placement Problem (DPP) is a critical hardware design optimization issue that involves finding the optimal placement of decoupling capacitors (decap) within a power distribution network (PDN) to enhance power integrity (PI) \citep{popovich2007power, shibasaka2013decoupling, erdin2019multi, park2023versatile}. Decoupling capacitors are hardware components that help reduce power noise and ensure a stable supply of power to the operating integrated circuits within hardware devices such as CPUs, GPUs, and AI accelerators \citep{popovich2007power}. The DPP is formulated as a black-box contextual optimization problem, where the goal is to determine the best positions for a set of decaps to maximize the PI objective. This objective is contextualized by the target hardware's feature vectors, with the constraint of a limited number of decaps. Interested readers can refer to \citep{kim2023devformer} for more details.



\paragraph{Instance generation.}
$10 \times 10$ PDN instances are used. We generate training and validation instances following \citet{kim2023devformer}. The test instances are directly drawn from \cite{kim2023devformer}.